\section{Comparison, composition and trend columns}
Available types are text, number, bar, bullet, interval, compare, dumbbell, stack
spark and heat. Compare and dumbbell cells use reference/current order; default labels
display current/reference, while delta displays current minus reference.
NA means both missing; 40/NA and NA/60 represent one missing endpoint.
\begin{demo}{Two views of the same comparison}
\begin{gradtable}[text width=22mm,graphic width=32mm,
  label width=12mm,header align=c,column sep=6pt,row stretch=1.8]
  \gradcolumn{name}{Method}
  \gradcolumn[type=compare,options={max=100,precision=0}]{bars}{Values}
  \gradcolumn[type=dumbbell,options={max=100,precision=0,
    compare label=delta,marker shape=diamond}]{change}{Change}
  \gradrow{A,40/68,40/68}
  \gradrow{B,75/60,75/60}
  \gradrow{C,50/NA,50/NA}
\end{gradtable}
\end{demo}
Stack and spark cells contain lists, so enclose each cell's list in braces.
Whole-cell NA is supported. Missing individual stack contributions are not zero
and must not be included in a total. A spark's missing last position retains a
missing label rather than borrowing an earlier observation.
\begin{demo}{Composition and a common-scale trend}
\begin{gradtable}[text width=22mm,graphic width=32mm,
  label width=12mm,header align=c,column sep=6pt,row stretch=1.8]
  \gradcolumn{name}{Project}
  \gradcolumn[type=stack,options={max=100,precision=0,
    stack colors={gradbarsBlue,gradbarsTeal,gold}}]{mix}{Composition}
  \gradcolumn[type=spark,options={spark range=fixed,
    min=0,max=100,precision=0}]{trend}{Five periods}
  \gradrow{A,{20,35,25},{35,48,NA,65,80}}
  \gradrow{B,{30,20,15},{80,70,65,NA,NA}}
  \gradrow{C,NA,NA}
\end{gradtable}
\end{demo}
Keep stack categories in a consistent order or register named categories.
Use a negative min for negative contributions. Use fixed spark ranges to compare
magnitudes; auto scales each cell separately and only supports shape comparison.
